% ===================================================================
%  TABLEAU N° 15 — Le Marché. — Les Comestibles.
%  Feuillets 84 a 87 du fac-simile = folios 82 a 85.
%  Correspondance : folio imprime = numero de feuillet - 2.
%
%  Le francais ne connait aucun des intertitres idists de ce tableau
%  (« Vendado e Komprado », « Chipa e Chera vari », « La Spicaro »,
%  « La Klientaro », « La Spektaji », « Atencez la policisto ! ») :
%  Rochelle laisse courir les 18 alineas sans subdivision apres le
%  titre. Seul « Le Marché. --- Les Comestibles. » repond a « La
%  Merkato. --- La Manjebli. » ; il porte donc la cle t15-tit-1 comme
%  dans texto/io, mais compose en \VUtitre (gras, bas de casse) et non
%  en \VUpk, conformement a ce que montre le fac-simile francais (meme
%  traitement que dans les tableaux 1 et 2).
%
%  Deux coquilles d'imprimerie relevees telles quelles (feuillet 85,
%  85) :
%   - « une dinde. une oie. un canard » : le compositeur a mis un
%     point la ou le sens appelle une virgule (repete deux fois) ;
%   - « côté d'elle, 'est très consciencieuse » : le « c » de « c'est »
%     manque, ne laissant que l'apostrophe.
%  Le renvoi des « petits pains » (feuillet 86) est imprime (32) alors
%  que la logique de la serie (cf. « semli » (82) en ido) voudrait
%  (82) ; le chiffre est net a l'agrandissement, donc conserve tel
%  quel plutot que corrige.
%
%  ATTENTION — LE FAC-SIMILE FRANCAIS EST UNE PRISE DE VUE, non un
%  passage au scanner : contraste inegal, quelques taches d'encre sur
%  les renvois chiffres. Chaque renvoi douteux a ete recoupe avec
%  texto/io/24-tabelo-15.tex, qui porte la meme legende chiffree, et
%  confirme un a un (halle=1 ... sac=105).
% ===================================================================

% --- Feuillet 84 (folio 82, ouverture de tableau) -------------------
\begin{VUpage}[84]{82}
%%K t15-apar-1 apar
\VUsaut{2.0mm}
\VUtitre{15.0pt}{60}{TABLEAU N\textsuperscript{o} 15}
\VUsaut{2.4mm}

%%K t15-tit-1 sub
\VUtitre{11.4pt}{0}{Le Marché. --- Les Comestibles.}
\VUsaut{4.0mm}

%%K t15-01-1 p
\VUblancAlinea
1. --- La plupart des commerçants qui nous fournis\cc
sent les denrées nécessaires à notre nourriture se\nl
trouvent réunis sous la \VUgras{halle} \textsuperscript{(1)} du \VUgras{marché couvert,}\nl
ou dans les rues voisines.

%%K t15-02-1 p
\VUblancAlinea
2. --- Le \VUgras{charcutier} \textsuperscript{(2)}, qui vend le \VUgras{jambon} \textsuperscript{(3)}, le\nl
\VUgras{boudin} \textsuperscript{(4)}, la \VUgras{saucisse} \textsuperscript{(5)}, l'\VUgras{andouillette} \textsuperscript{(6)} et le\nl
\VUgras{lard}\textsuperscript{(7)}, a pour voisine la \VUgras{marchande de beurre}\nl
\VUgras{et de fromage} \textsuperscript{(8)} dont l'étalage est garni de \VUgras{fro}\cc
\VUgras{mages de Hollande} \textsuperscript{(9)} à la croûte rouge, de\nl
\VUgras{camemberts} \textsuperscript{(10)}, d'énormes \VUgras{pains de gruyère} \textsuperscript{(11)}\nl
et de \VUgras{mottes de beurre} \textsuperscript{(12)} frais.

%%K t15-03-1 p
\VUblancAlinea
3.--- Devant elle, une \VUgras{marchande de légumes} \textsuperscript{(13)}\nl
offre à ses clientes de belles \VUgras{tomates} \textsuperscript{(14)}, des \VUgras{choux}\cc
\VUgras{fleurs} \textsuperscript{(15)}, de la \VUgras{salade} \textsuperscript{(16)}, des \VUgras{choux pommés} \textsuperscript{(17)},\nl
des \VUgras{citrouilles} \textsuperscript{(19)}, des \VUgras{melons} \textsuperscript{(18)} et même des\nl
\VUgras{champignons} \textsuperscript{(20)}. Mais un \VUgras{garçon brasseur}, qui a\nl
arrêté sa \VUgras{voiture de livraison} \textsuperscript{(21)} chargée de \VUgras{ton}\cc
\VUgras{neaux de bière} \textsuperscript{(22)} juste devant son étalage, empêche\nl
ses pratiques d'approcher. Une jardinière lui apporte\nl
une corbeille d'\VUgras{artichauts} \textsuperscript{(23)}.

%%K t15-04-1 p
\VUblancAlinea
4. --- Sous le marché, l'animation est grande, car\nl
c'est l'heure de la \VUgras{criée} \textsuperscript{(24)}. Les \VUgras{ménagères}, \textsuperscript{(26)} qui\nl
viennent d'y faire leurs achats, s'arrêtent en passant\nl
devant la boutique de la \VUgras{tripière} \textsuperscript{(25)}, pour acheter du\nl
\VUgras{gras-double} ou de la \VUgras{tête de veau.}

%%K t15-05-1 p
\VUblancAlinea
5. --- Un gros \VUgras{cuisinier} \textsuperscript{(27)} marchande un très beau\nl
\VUgras{thon} à la \VUgras{marchande de poissons} \textsuperscript{(28)}, qui exagère\nl
ses prix à plaisir. Elle a un superbe arrivage d'\VUgras{an}\cc
\VUgras{guilles,} de \VUgras{soles,} de \VUgras{truites} et de \VUgras{carpes.} Sa\nl
\VUgras{morue} et ses \VUgras{moules} sont très fraîches.

%%K t15-06-1 p
\VUblancAlinea
6. --- La \VUgras{marchande de volaille} \textsuperscript{(29)}, installée à\nl
côté d'elle, 'est très consciencieuse. Lorsqu'elle vend
\parplein
\end{VUpage}

% --- Feuillet 85 (folio 83) ------------------------------------------
\begin{VUpage}[85]{83}
%%K t15-06-1 p suite
\VUcontinue
une \VUgras{dinde}. une \VUgras{oie}. un \VUgras{canard} ou un simple \VUgras{poulet,}\nl
elle ne demande que le plus juste prix.

%%K t15-07-1 p
\VUblancAlinea
7. --- A l'angle de la place, au premier étage, est\nl
établi un \VUgras{marchand de toiles} \textsuperscript{(30)} chez lequel les\nl
acheteuses sont toujours sûres de trouver un très joli\nl
assortiment de \VUgras{nappes}, de \VUgras{serviettes}, \VUgras{tabliers} ou\nl
\VUgras{torchons.} Au même étage, un \VUgras{coutelier} \textsuperscript{(31)} a ins\cc
tallé son \VUgras{atelier.} Il repasse les \VUgras{couteaux} des mar\cc
chandes de la halle et fabrique pour les jardiniers des\nl
\VUgras{serpes} en acier de première qualité.

%%K t15-08-1 p
\VUblancAlinea
8. --- Au-dessous de son atelier s'est ouvert depuis\nl
peu de temps un important magasin de comestibles.\nl
Ce n'est plus simplement la modeste \VUgras{épicerie} \textsuperscript{(32)}\nl
d'autrefois, qui se contentait de vendre les choses\nl
usuelles, le \VUgras{thé} \textsuperscript{(33)}, le \VUgras{chocolat} \textsuperscript{(34)}, le \VUgras{sucre en}\nl
\VUgras{pain} \textsuperscript{(37)} et le \VUgras{savon} \textsuperscript{(38)}. Ici on vend de tout, des\nl
\VUgras{gâteaux secs}, des \VUgras{conserves} \textsuperscript{(35)}, des \VUgras{liqueurs} \textsuperscript{(36)},\nl
\VUgras{rhum, cognac, kirsch, anisette} et \VUgras{curaçao.} On\nl
y trouve aussi bien du \VUgras{gibier} : \VUgras{lièvre} \textsuperscript{(39)}, \VUgras{grives} \textsuperscript{(40)},\nl
\VUgras{perdrix} \textsuperscript{(41)}, \VUgras{cailles} \textsuperscript{(42)}, \VUgras{canard sauvage} \textsuperscript{(43)} ou\nl
\VUgras{faisan} \textsuperscript{(44)}, que des fruits exotiques tels que les \VUgras{bana}\cc
\VUgras{nes} \textsuperscript{(46)} et les \VUgras{ananas.}

%%K t15-09-1 p
\VUblancAlinea
9. --- Un \VUgras{garçon} \textsuperscript{(45)} épicier, très complaisant, est\nl
prêt à servir ce que l'on désire : \VUgras{confiture} \textsuperscript{(47)} (de\nl
\VUgras{groseille} ou de \VUgras{coing}), \VUgras{graisse} \textsuperscript{(48)}, \VUgras{cornichons} \textsuperscript{(49)}\nl
ou bien \VUgras{sardines} \textsuperscript{(50)} salées.

%%K t15-09-2 p
C'est fort commode pour les \VUgras{clientes} \textsuperscript{(51)}, qui trou\cc
vent, à côté des choses les plus vulgaires, les \VUgras{pri}\cc
\VUgras{meurs} les plus rares et les \VUgras{fruits} les plus savou\cc
reux : \VUgras{poires} \textsuperscript{(52)} fondantes, \VUgras{prunes} \textsuperscript{(53)}, \VUgras{raisins} \textsuperscript{(54)},\nl
\VUgras{pêches} \textsuperscript{(55)}, \VUgras{amandes} \textsuperscript{(56)} et \VUgras{noix} \textsuperscript{(57)}.

%%K t15-10-1 p
\VUblancAlinea
10. --- Le \VUgras{riz} \textsuperscript{(58)} et les \VUgras{lentilles} \textsuperscript{(59)} sont à côté de\nl
la caisse de \VUgras{harengs-saurs} \textsuperscript{(60)}; la \VUgras{pomme de}\nl
\VUgras{terre} \textsuperscript{(61)} et le \VUgras{haricot} \textsuperscript{(63)} voisinent avec les \VUgras{pom}\cc
\VUgras{mes} \textsuperscript{(62)}, les \VUgras{figues} \textsuperscript{(64)}, les \VUgras{pruneaux} \textsuperscript{(65)} et les\nl
\VUgras{noisettes} \textsuperscript{(66)}. On trouve dans ce magasin des \VUgras{œufs} \textsuperscript{(67)}
\parplein
\end{VUpage}

% --- Feuillet 86 (folio 84) ------------------------------------------
\begin{VUpage}[86]{84}
%%K t15-10-1 p suite
\VUcontinue
frais et des \VUgras{olives} \textsuperscript{(68)}; aussi \VUgras{paniers} \textsuperscript{(70)} et \VUgras{filets à}\nl
\VUgras{provisions} \textsuperscript{(73)} s'emplissent-ils à vue d'œil.

%%K t15-11-1 p
\VUblancAlinea
11. --- Le \VUgras{café} \textsuperscript{(72)} de cette excellente maison a\nl
acquis parmi les \VUgras{bonnes} \textsuperscript{(69)} et les cuisinières une\nl
véritable réputation, car c'est le vieil \VUgras{épicier} \textsuperscript{(71)} qui\nl
le grille lui-même avec le plus grand soin.

%%K t15-12-1 p
\VUblancAlinea
12. --- Tout le monde ne vient pas au marché pour\nl
y faire des provisions. Dans le pittoresque va-et-vient\nl
de la petite rue qui conduit à la halle, on voit les\nl
personnages les plus différents. A côté d'un brave\nl
\VUgras{garçon d'abattoir} \textsuperscript{(75)}, qui pousse devant lui\nl
une \VUgras{brouette} \textsuperscript{(74)}, voici deux soldats en permission\nl
tout fiers d'exhiber leur brillant uniforme. Le \VUgras{hus}\cc
\VUgras{sard} \textsuperscript{(76)}, en \VUgras{dolman} bleu et \VUgras{shako} à plumes, et le\nl
\VUgras{caporal de zouaves,} en culotte bouffante et coiffé\nl
d'une chéchia.

%%K t15-13-1 p
\VUblancAlinea
13. --- Le \VUgras{boulanger} \textsuperscript{(80)}, qui se tient sur le seuil\nl
de la \VUgras{boulangerie} \textsuperscript{(78)}, vient de pétrir la \VUgras{pâte} de son\nl
\VUgras{pain} \textsuperscript{(79)} et de retirer du four les \VUgras{croissants} \textsuperscript{(81)}, les\nl
\VUgras{petits pains} \textsuperscript{(32)} et les \VUgras{miches} \textsuperscript{(83)} qui sont dans la\nl
vitrine.

%%K t15-14-1 p
\VUblancAlinea
14. --- Au-dessus de la boulangerie, habite une\nl
adroite \VUgras{lingère} \textsuperscript{(84)} qui emploie plusieurs ouvrières\nl
\VUgras{brodeuses} \textsuperscript{(85)}. Auprès d'elle demeure une \VUgras{blanchis}\cc
\VUgras{seuse-repasseuse} \textsuperscript{(86)}, qui amidonne le \VUgras{linge} \textsuperscript{(88)}\nl
après l'avoir lavé et fait sécher, et qui le repasse\nl
avec un \VUgras{fer} \textsuperscript{(87)}.

%%K t15-15-1 p
\VUblancAlinea
15. --- Dans la \VUgras{boucherie} \textsuperscript{(89)}, située au rez-de-\nl
chaussée, on vend toutes sortes de \VUgras{viandes}, \VUgras{mou}\cc
\VUgras{ton} \textsuperscript{(90)}, \VUgras{bœuf} \textsuperscript{(94)} ou \VUgras{veau} \textsuperscript{(92)}, sous forme de \VUgras{gigots,}\nl
\VUgras{biftecks, rôtis} et \VUgras{côtelettes.} Le \VUgras{garçon bou}\cc
\VUgras{cher} \textsuperscript{(93)} détaille toutes ces viandes et met de côté\nl
les \VUgras{os} \textsuperscript{(96)} à moelle. A la porte de la boucherie se\nl
tient une \VUgras{écaillère} \textsuperscript{(97)} qui vend des \VUgras{huîtres} \textsuperscript{(98)}.\nl
Elle les ouvre, si on le désire.

%%K t15-16-1 p
\VUblancAlinea
16. --- Tous ces commerçants paient une patente\nl
ou un droit de plaçage; c'est pourquoi l'\VUgras{agent de}
\parplein
\end{VUpage}

% --- Feuillet 87 (folio 85) ------------------------------------------
\begin{VUpage}[87]{85}
%%K t15-16-1 p suite
\VUcontinue
\VUgras{police} \textsuperscript{(100)} pourchasse impitoyablement les malheu\cc
reuses \VUgras{marchandes des quatre saisons} \textsuperscript{(99)} qui\nl
leur font concurrence en transportant dans leurs\nl
voitures des \VUgras{carottes}, des \VUgras{radis}, des \VUgras{navets,}\nl
des \VUgras{poireaux} ou bien des \VUgras{bottes d'asperges}, des\nl
\VUgras{pois verts}, des \VUgras{épinards}, de l'\VUgras{oseille}, du \VUgras{cresson}\nl
et du \VUgras{persil.}

%%K t15-17-1 p
\VUblancAlinea
17. --- Le petit garçon qui vend de l'\VUgras{ail} \textsuperscript{(101)} et des\nl
\VUgras{oignons} sait bien que lui non plus ne doit pas ven\cc
dre sa marchandise auprès du marché. Il s'enfuit en\nl
courant, au risque de renverser le panier de \VUgras{crabes,}\nl
d'\VUgras{écrevisses} et de \VUgras{crevettes} du \VUgras{marchand} \textsuperscript{(102)}\nl
de \VUgras{langoustes} et de \VUgras{homards.}

%%K t15-18-1 p
\VUblancAlinea
18. --- Tout ce monde s'agite et fait un bruit\nl
effroyable, mais cela n'empêche pas d'entendre dis\cc
tinctement l'appel du \VUgras{vitrier} \textsuperscript{(103)} ambulant, ni le cri\nl
du \VUgras{charbonnier} \textsuperscript{(104)} qui vient de quitter un moment\nl
sa petite charrette pour livrer un \VUgras{sac} \textsuperscript{(105)} de \VUgras{charbon.}

\VUsaut{6.0mm}
\VUfilet{22mm}
\end{VUpage}
